\subsection{ORDER FUNCTION}

Let \(G\) be a group with identity element \(e\) . Let \((\mathbf{G}_{\alpha})\) be a real filtration on \(G\) . For all \(x\) in \(G\) let \(I_x\) denote the set of real numbers \(\alpha\) such that \(x \in G_{\alpha}\) . If \(\alpha \in I_x\) and \(\beta < \alpha\) , then \(\beta \in I_x\) and hence \(I_x\) is an interval (General Topology, Chapter IV, § 2, no. 4, Proposition 1). Using relation (1), we see that \(I_x\) contains its least upper bound when this is finite. Therefore \(I_x\) is of the form \(] - \infty\) , \(v(x)] \cap R\) with \(v(x) \in \overline{\mathbf{R}}\) ; we have \(v(x) = \sup(\alpha \mid x \in G_{\alpha})\) .

The mapping \(v\) of \(G\) into \(\overline{\mathbf{R}}\) is called the order function associated with the real filtration \((G_{\alpha})\) and \(v(x)\) is called the order of \(x\) . This mapping has the following properties:

\begin{enumerate}
\def\labelenumi{(\alph{enumi})}
\item
  For \(x \in \mathbf{G}\) and \(\alpha \in \mathbf{R}\) , the relations \(x \in \mathbf{G}_{\alpha}\) and \(v(x) \geqslant \alpha\) are equivalent.
\item
  For x, y in G,
\end{enumerate}

\begin{enumerate}
\def\labelenumi{(\arabic{enumi})}
\setcounter{enumi}{1}
\tightlist
\item
\end{enumerate}

\[
v (x ^ {- 1}) = v (x), \quad v (e) = + \infty .\tag{2}
\]

\[
v (x y) \geqslant \inf (v (x), v (y)).\tag{3}
\]

Further, we have equality in (3) if \(v(x) > v(y)\) .

\begin{enumerate}
\def\labelenumi{(\alph{enumi})}
\setcounter{enumi}{2}
\tightlist
\item
  For all \(\alpha \in \mathbf{R}\) , let \(\mathrm{G}_{\alpha}^{+}\) denote the set of \(x \in \mathrm{G}\) such that \(v(x) > \alpha\) . Then \(\mathrm{G}_{\alpha}^{+} = \bigcup_{\beta > \alpha} \mathrm{G}_{\beta}\) and in particular \(\mathrm{G}_{\alpha}^{+}\) is a subgroup of \(\mathrm{G}\) .
\end{enumerate}

Conversely, let \(v\) be a mapping of \(G\) into \(\overline{\mathbf{R}}\) satisfying relations (2) and (3). For all \(\alpha \in \mathbf{R}\) , let \(G_{\alpha}\) be the set of \(x \in G\) such that \(v(x) \geqslant \alpha\) . Then \((G_{\alpha})_{\alpha \in \mathbf{R}}\) is a real filtration of \(G\) and \(v\) is the order function associated with this filtration.

For the filtration \((\mathbf{G}_{\alpha})\) to be integral, it is necessary and sufficient that \(v\) map \(\mathbf{G}\) into \(\mathbf{Z} \cup \{+\infty, -\infty\}\) . For it to be exhaustive (resp. separated), it is necessary and sufficient that \(v^{-1}(-\infty) = \varnothing\) (resp. \(v^{-1}(+\infty) = \{e\}\) ).

Let A be a K-algebra (resp. unital K-algebra). By the above, the relation

\[
" x \in \mathrm{A} _ {\alpha} \Leftrightarrow v (x) \geqslant \alpha \quad \text { for } x \in \mathrm{A} \text { and } \alpha \in \mathbf {R}"
\]

defines a bijection of the set of exhaustive real filtrations \((\mathbf{A}_{\alpha})_{\alpha\in\mathbf{R}}\) compatible with the algebra (resp. unital algebra) structure on A onto the set of mappings \(v: A \to \overline{\mathbf{R}}\) not taking the value \(-\infty\) and satisfying axioms (4) to (7) (resp. (4) to (8)) below:

\begin{enumerate}
\def\labelenumi{(\arabic{enumi})}
\setcounter{enumi}{3}
\tightlist
\item
\end{enumerate}

\[
v (x + y) \geqslant \inf (v (x), v (y)) \quad (x, y \text {   in   A })\tag{5}
\]

\[
v (- x) = v (x) \quad (x \in \mathrm{A})\tag{6}
\]

\[
v (\lambda x) \geqslant v (x) \quad (\lambda \in \mathrm{K}, x \in \mathrm{A})\tag{7}
\]

\[
v (x y) \geqslant v (x) + v (y) \quad (x, y \text {   in   A })\tag{8}
\]

\[
v (1) \geqslant 0.
\]

Remark. If \(v(x)\) is not everywhere equal to \(+\infty\) , conditions (7) and (8) imply \(v(1) = 0\) .

